\documentclass{article}
\usepackage{decomp}
\begin{document}

\sourcefile{elife-08206-figures-v3.pdf}
\decompnote{The eLife ``Figures and figure supplements'' PDF of the article (13 pages): every main figure again with its legend, and the five figure supplements with their legends (Figure 1—figure supplements 1 and 2, Figure 2—figure supplement 1, Figure 3—figure supplement 1, Figure 4—figure supplement 1). The running header (eLife logo, ``Research Article'', ``Neuroscience'') and footer (``Pallotto et al. eLife 2015;4:e08206. DOI: 10.7554/eLife.08206'', ``n of 13'') are left out. The images of the main figures are identical to those in the article PDF (pixel for pixel), so they reuse \texttt{figures/fig1.png} \ldots\ \texttt{fig5.png}; their descriptions and data are in \texttt{main.tex}. The figure-supplement images were rendered from this PDF; the full-size eLife JPEGs in \texttt{papers/pallotto2015/elife\_figure\_supplements/} have the same pixel dimensions (mean difference 0.4--3 grey levels, JPEG re-encoding) and were used for the data extraction. As in the article PDF, the text layer renders μ as ``m''; the page images show μm, transcribed here.}

\pdfpage{1}
\noindent RESEARCH ARTICLE\par

\section*{Figures and figure supplements}

\noindent Extracellular space preservation aids the connectomic analysis of neural circuits

\noindent \textbf{Marta Pallotto \emph{et al}}

\pdfpage{2}
\begin{center}
\includegraphics[width=0.95\linewidth,height=0.55\textheight,keepaspectratio]{figures/fig1.png}
\end{center}
\noindent\textbf{Figure 1.} ECS preservation in acutely isolated tissues. (\textbf{A}) A model of the changes brain tissue undergoes during chemical tissue fixation and possible interventions to preserve ECS. (\textbf{B}) ECS fraction of mouse olfactory bulb (x), retina (o), and cortex (□) correlates with increasing fixative vehicle concentration. Dashed line is the linear fit to olfactory bulb data, solid line is the linear fit to cortex data (olfactory bulb: R = 0.80, p = 0.0006; cortex: R = 0.79, p = 0.0113). (\textbf{C–F}) EM images from the external plexiform layer of the mouse olfactory bulb from a perfused brain (\textbf{C}) or from acute sections fixed with increasing buffer concentrations (\textbf{D–F}). ECS fractions are 0.6\% (\textbf{C}), 5.8\% (\textbf{D}), 11.3\% (\textbf{E}), and 23.9\% (\textbf{F}). Scale bars: 1 μm (\textbf{C–F}). ECS: Extracellular space; EM: Electron microscopy.

\noindent DOI: 10.7554/eLife.08206.003
\alttext{Same image as Figure 1 of the article; see the description and \texttt{figures/fig1b.csv} in \texttt{main.tex}.}

\pdfpage{3}
\begin{center}
\includegraphics[width=0.95\linewidth,height=0.55\textheight,keepaspectratio]{figures/sfig1-1.png}
\end{center}
\noindent\textbf{Figure 1—figure supplement 1.} Change of tissue dimensions and uniformity of ECS preserved sections. (\textbf{A}) Vibratome sections were cut from the cerebral cortex of acutely isolated brains (for ECS preservation) or from perfused animals. Section thickness was nominally 500 μm as determined by the slice thickness setting of the vibratome. Sections were stained for EM and the cross-sectional thickness measured from EM images (mean +/- SEM; n = 20 slices from 6 animals, one-way ANOVA: p = 0.46). (\textbf{B}) Sample images from the edges (\textbf{Bi}, \textbf{Biii}) and center (\textbf{Bii}) of a 500-μm-thick section of the mouse olfactory bulb. ECS is evenly preserved across the section. ECS: Extracellular space; EM: Electron microscopy.

\noindent DOI: 10.7554/eLife.08206.004
\alttext{(A, top left) Point plot with error bars (filled circles, mean +/- SEM): y ``Slice thickness (um)'', linear 0–600, ticks every 100; x ``Fix vehicle osmolarity (mOsm)'' with ticks 200, 300, 400, 500 and a category ``perfused'' left of an axis break (//). Points: perfused ≈ 489 μm (bar ≈ 475–504); ≈ 200 mOsm ≈ 509 (≈ 487–532); ≈ 302 mOsm ≈ 504 (≈ 496–513); ≈ 372 mOsm ≈ 527 (≈ 514–540); ≈ 407 mOsm ≈ 494 with no bar visible outside the marker. All lie close to the nominal 500 μm. (B, top right) Schematic: a tall black-outlined rectangle (the section in cross-section) with three dashed corner-bracket boxes side by side at mid-height, labelled i, ii, iii (left edge, centre, right edge). (Bi, Bii, Biii, bottom row) Three square greyscale EM images of olfactory-bulb neuropil from those positions. All three show similarly preserved extracellular space: light (white) irregular spaces between grey cellular profiles, dark mitochondria; Bii contains a dark myelinated axon cross-section at the upper right. No scale bar is printed.}
\figuredata{A}{figures/sfig1-1a.csv}{figures/sfig1-1a.py}{tier 3, digitized from the eLife JPEG \texttt{elife-08206-fig1-figsupp1-v3.jpg} (marker cores and error-bar ends)}

\pdfpage{4}
\begin{center}
\includegraphics[width=0.95\linewidth,height=0.55\textheight,keepaspectratio]{figures/sfig1-2.png}
\end{center}
\noindent\textbf{Figure 1—figure supplement 2.} Effect of temperature on the preservation of ECS. (\textbf{A},\textbf{B}) Vibratome sections were cut from the olfactory bulbs of acutely isolated brains at either 4°C or 25°C and then fixed at 25°C in 0.1 M CB + 2\% GA. (\textbf{C},\textbf{D}) Vibratome sections were cut from the olfactory bulbs of acutely isolated brains at either 4°C or 25°C and then fixed at 25°C in 0.2 M CB + 2\% GA. ECS preservation was similar under both conditions indicating the temperature of solutions during tissue dissection is not a major determinant for ECS preservation. Scale bars: 1 μm. ECS: Extracellular space; EM: Electron microscopy.

\noindent DOI: 10.7554/eLife.08206.005
\alttext{Four square greyscale EM images of olfactory-bulb neuropil in a 2 × 2 grid, each with a 1 μm scale bar at the lower right and a title: (A) ``25°C dissection, 25°C fixation (0.1 M CB)''; (B) ``4°C dissection, 25°C fixation (0.1 M CB)''; (C) ``25°C dissection, 25°C fixation (0.2 M CB)''; (D) ``4°C dissection, 25°C fixation (0.2 M CB)''. A and B (0.1 M CB): rounded, closely packed profiles with thin membranes and little visible extracellular space. C and D (0.2 M CB): many white extracellular gaps between smaller profiles, more in D; D also shows a large smooth (nuclear or somatic) region at the upper right. Within each buffer pair the two dissection temperatures look alike.}
\decompnote{The panel titles and this legend give 25°C dissection and fixation temperatures; the main text says sections were cut ``at 4°C or 20°C'' and fixed at 20°C.}

\pdfpage{5}
\begin{center}
\includegraphics[width=0.95\linewidth,height=0.55\textheight,keepaspectratio]{figures/fig2.png}
\end{center}
\noindent\textbf{Figure 2.} Automated 2D segmentation of extracellular space preserved data. (\textbf{A}) Four raw EM images (first column) of varying ECS fractions (0.6, 5.8, 11.3, and 23.9\%) were analyzed. The lowest ECS fraction (0.6\%) is the perfused preparation while the others are the acute slice preparations. Pixels in each image were annotated (second column) as either intracellular (white), extracellular (gray), or plasma membrane (black). Representative test image pixel classifications (third column) yielded network classification errors of 9.3, 8.1, 6.7, and 7.2\% for the four examples shown. Intracellular pixel probabilities were thresholded and segmented (fourth column) to yield minimum Rand errors of 0.11, 0.06, 0.06, and 0.07 for the four examples shown. (\textbf{B}) Pixel classification errors plotted versus ECS fraction from an ninefold cross-validation, median values in red, shows no significant difference between perfused and acute preparations (K-W test: p=0.05). (\textbf{C}) Rand error plotted versus ECS fraction for the ninefold cross-validation, median values in red, shows a significant difference between perfused and acute preparations (K-W test: p<0.01; Wilcoxon rank-sum test: p=0.0017 [5.8\%], p=0.0091 [11.3\%], p=0.031 [23.9\%]). (\textbf{D}) Total number of splits and merger per object in each test image plotted versus ECS fraction for the ninefold cross-validation, median values in red, shows a significant difference between perfused and acute preparations (K-W test: p<0.001; Wilcoxon rank-sum test: p=0.00049 [5.8\%], p=0.00035 [11.3\%], p=0.00035 [23.9\%]). Scale bars: 2 μm. p-Values: *p<0.05, **p<0.01, ***p<0.001. ECS: Extracellular space; EM: Electron microscopy.

\noindent \emph{Figure 2 continued on next page}

\pdfpage{6}
\noindent \emph{Figure 2 continued}

\noindent DOI: 10.7554/eLife.08206.006

\noindent The following source data is available for figure 2:

\noindent\textbf{Source data 1.} Raw image and expert labels of 0.6\% ECS data.

\noindent DOI: 10.7554/eLife.08206.007

\noindent\textbf{Source data 2.} Raw image and expert labels of 5.8\% ECS data.

\noindent DOI: 10.7554/eLife.08206.008

\noindent\textbf{Source data 3.} Raw image and expert labels of 11.3\% ECS data.

\noindent DOI: 10.7554/eLife.08206.009

\noindent\textbf{Source data 4.} Raw image and expert labels of 23.9\% ECS data.

\noindent DOI: 10.7554/eLife.08206.010
\alttext{Same image as Figure 2 of the article; see the description and \texttt{figures/fig2b.csv}, \texttt{fig2c.csv}, \texttt{fig2d.csv} in \texttt{main.tex}.}

\pdfpage{7}
\begin{center}
\includegraphics[width=0.95\linewidth,height=0.55\textheight,keepaspectratio]{figures/sfig2-1.png}
\end{center}
\noindent\textbf{Figure 2—figure supplement 1.} Warping error of automated segmentations. (\textbf{A}) Split/merger curves parameterized by threshold of the network output. Shaded area is the standard error of the mean, lines are the mean number of splits/mergers per object, and the circles are the mean of the minimum total splits + mergers per object. (\textbf{B}) Network segmentations at the minimum warping error threshold, same regions as in Figure 2A. Warping errors were 0.68 (0.6\% ECS), 0.32 (5.8\% ECS), 0.52 (11.3\% ECS), and 0.56 (23.9\% ECS) for the four examples shown. (\textbf{C}) Cumulative distributions of the number of ground truth intracellular pixels per object for the four datasets. The difference in these distributions motivated the use of a Bernoulli sampling procedure for calculating the adapted Rand Error metrics in Figure 2C. Scale bars: 2 μm. ECS: Extracellular space.

\noindent DOI: 10.7554/eLife.08206.011
\alttext{(A, top left) Line plot with shaded bands: y ``Mergers / Object'' 0–0.6, x ``Splits / Object'' 0–0.6 (ticks every 0.1). Legend: ``0.6\% ECS'' (purple), ``5.8\% ECS'' (olive green), ``11.3\% ECS'' (orange), ``23.9\% ECS'' (dark red). The 0.6 \% curve enters at the top (y = 0.6) near x ≈ 0.26, falls through its circle at ≈ (0.32, 0.45) and reaches ≈ 0.30 at x = 0.6, inside a wide purple band (about ±0.04). The three other curves start near (0.12–0.13, 0.30), have their circles at ≈ (0.15, 0.25) (11.3 \%), (0.17, 0.22) (23.9 \%) and (0.19, 0.20) (5.8 \%), and flatten to ≈ 0.10–0.12 at x = 0.6; they and their overlapping bands lie on top of each other, with 11.3 \% slightly highest at large x. (B, right) Two columns × four rows of square images, column titles ``expert labeled'' and ``segmented network output at min (Splits + Mergers) / Object''; row labels ``perfused (0.6 \% ECS)'', ``acute slice (5.8 \% ECS)'', ``acute slice (11.3 \% ECS)'', ``acute slice (23.9 \% ECS)''. Expert labels: white objects, black membranes, grey extracellular space (almost none in the perfused row, large grey areas in the 23.9 \% row); a 2 μm scale bar under each. Segmentations: objects in four colours (blue, green, yellow, red) on black. (C, bottom left) Cumulative distribution plot: y ``Cumulative Distribution Function (\%)'' 0–100, x ``2D GT Pixels / Object'' logarithmic 10–10\textsuperscript{5}. Four S-shaped curves in the legend colours of A, all 0 below ≈ 18 pixels and reaching 100 \% near 3–8 × 10\textsuperscript{4}. Median crosses (×) at 50 \%: 23.9 \% at ≈ 280 pixels, 11.3 \% and 5.8 \% (overlapping) at ≈ 470, 0.6 \% at ≈ 750; the 23.9 \% curve lies left of (smaller objects than) the others over most of the range.}
\figuredata{A}{figures/sfig2-1a.csv}{figures/sfig2-1.py}{tier 3, digitized from the eLife JPEG (colour masks per pixel column; band edges for the 0.6 \% curve only; parts of overlapping curves hidden)}
\figuredata{C}{figures/sfig2-1c.csv}{figures/sfig2-1.py}{tier 3, digitized from the eLife JPEG; a distribution recomputed from the Figure 2 source data (expert labels) is in \texttt{figures/sfig2-1c\_sourcedata.csv} (same medians within about 5 \%, but the plotted curves exclude objects below ≈ 18 pixels by an unstated procedure)}

\pdfpage{8}
\begin{center}
\includegraphics[width=0.95\linewidth,height=0.55\textheight,keepaspectratio]{figures/fig3.png}
\end{center}
\noindent\textbf{Figure 3.} Automated 3D segmentation of extracellular space preserved data. (\textbf{A}) A 10 x 10 x 12 μm\textsuperscript{3} 3D SBEM cube (left) of perfused tissue (LowECS), a dense skeletonization of all neurites in the volume (middle), and an automated segmentation of the volume (right). (\textbf{B}) A 10 x 10 x 12 μm\textsuperscript{3} 3D SBEM cube (left) of ECS-preserved tissue (HighECS), a dense skeletonization of all neurites in the volume (middle), and an automated segmentation of the volume (right). (\textbf{C}) Adapted Rand error of automated segmentations compared to dense skeletonizations as a function of the segmentation threshold for the HighECS (brown) and LowECS (blue) data volumes. (\textbf{D}) The fraction of merged skeleton nodes versus split skeleton edges as a function of network threshold. Solid points indicate the minimum sum of the error fractions. (\textbf{E}) The total error-free skeleton path length recovered as a function of network threshold. Shaded regions in \textbf{C–E} are the 99\% confidence intervals based on a Bernoulli sampling of the data. Scale bars in \textbf{A}, \textbf{B} = 1 μm. SBEM: Serial block-face scanning electron microscopy.

\noindent DOI: 10.7554/eLife.08206.012

\noindent The following source data is available for figure 3:

\noindent\textbf{Source data 1.} M0027\_11 dense skeletonization, Knossos NML file.

\noindent DOI: 10.7554/eLife.08206.013

\noindent\textbf{Source data 2.} M0007\_33 dense skeletonization, Knossos NML file.

\noindent DOI: 10.7554/eLife.08206.014

\noindent\textbf{Source data 3.} M0007\_33 training data cubes.

\noindent \emph{Figure 3 continued on next page}

\pdfpage{9}
\noindent \emph{Figure 3 continued}

\noindent DOI: 10.7554/eLife.08206.015

\noindent\textbf{Source data 4.} M0027\_11 training data cubes.

\noindent DOI: 10.7554/eLife.08206.016
\alttext{Same image as Figure 3 of the article; see the description and \texttt{figures/fig3c.csv}, \texttt{fig3d.csv}, \texttt{fig3e.csv} in \texttt{main.tex}.}

\pdfpage{10}
\begin{center}
\includegraphics[width=0.95\linewidth,height=0.55\textheight,keepaspectratio]{figures/sfig3-1.png}
\end{center}
\noindent\textbf{Figure 3—figure supplement 1.} Cumulative distributions for 3D hand-labeled volumes and skeletonizations of the LowECS (blue) and HighECS (brown) data volumes. (\textbf{A}) Cumulative distribution frequency (CDF) of number of intracellular ground truth voxels per object. The distribution medians are significantly different (Wilcoxon rank-sum test, p <0.0001). (\textbf{B}) CDF of path length per skeleton, medians 10.2 (LowECS) and 10.1 (HighECS) μm. The distributions are not significantly different (two-sample Kolmogorov–Smirnov (KS) test, p = 0.91). (\textbf{C}) CDF of the number of nodes annotated per skeleton, medians 16 (LowECS) and 17 (HighECS) nodes. The distributions not significantly different (two-sample KS test, p = 0.77).

\noindent DOI: 10.7554/eLife.08206.017
\alttext{Three cumulative distribution plots, each with y ``Cumulative Distribution Function (\%)'' 0–100 (ticks 0, 25, 50, 75, 100), LowECS in blue and HighECS in orange-brown, and a × of the series colour at each median (50 \%). (A, top left) x ``3D GT Voxels / Object'', logarithmic 10 to ≈ 2 × 10\textsuperscript{6}; legend ``LowECS'', ``HighECS''. HighECS rises from ≈ 25 voxels and lies left of LowECS throughout: medians ≈ 3.7 × 10\textsuperscript{3} (HighECS) and ≈ 1.2 × 10\textsuperscript{4} (LowECS) voxels; both reach 100 \% near 5–8 × 10\textsuperscript{5}. The HighECS curve has a shoulder near 19 \% at a few hundred voxels. (B, top right) x ``Path Length / Skeleton (um)'', logarithmic 0.1 to ≈ 160. The two step curves almost coincide: a few per cent below 2 μm, a steep rise between ≈ 3 and 20 μm, the medians (one × each, overlapping) at ≈ 10 μm, 100 \% at ≈ 60–100 μm. (C, bottom right) x ``Nodes / Skeleton'', linear 0–150. Both curves rise steeply from 0, the medians (overlapping ×) at ≈ 16–17 nodes, ≈ 90 \% at ≈ 45–55 nodes and 100 \% near 130–150 nodes; LowECS is slightly higher (more small skeletons) between ≈ 20 and 60 nodes. The lower left quadrant of the figure is empty.}
\figuredata{A}{figures/sfig3-1a.csv}{figures/sfig3-1.py}{tier 1, computed from the eLife source data (Figure 3—source data 3 and 4, training cubes): voxel count of every labelled intracellular object; checked against the printed curves}
\figuredata{B}{figures/sfig3-1b.csv}{figures/sfig3-1.py}{tier 3, digitized from the eLife JPEG; skeleton path lengths recomputed from the NML source data (Figure 3—source data 1 and 2) in \texttt{figures/sfig3-1b\_sourcedata.csv} match LowECS but not the plotted HighECS curve}
\figuredata{C}{figures/sfig3-1c.csv}{figures/sfig3-1.py}{tier 3, digitized from the eLife JPEG; node counts per skeleton from the NML files in \texttt{figures/sfig3-1c\_sourcedata.csv} (same caveat as B)}
\decompnote{``The distributions not significantly different'' is as printed.}

\pdfpage{11}
\begin{center}
\includegraphics[width=0.95\linewidth,height=0.55\textheight,keepaspectratio]{figures/fig4.png}
\end{center}
\noindent\textbf{Figure 4.} Extracellular space aids the identification of gap junctions. (\textbf{A}) Reconstructed dendrites of an AII amacrine cell (gray) and surface renderings of tight (yellow) and cleft (black) contacts. (\textbf{B-D}) Example EM images and volumetric reconstructions of (\textbf{B}) an AII (gray) to AII (blue) tight contact, (\textbf{C}) an AII (gray) to cone bipolar cell (green) tight contact (\textbf{C}), and (\textbf{D}) a cleft contact identified as a chemical synapse between a presynaptic amacrine cell (cyan) and the AII cell (gray). (\textbf{E}) Tight contact surface renderings color-coded by the identity of the synaptic partner: AII amacrine cell (blue), cone bipolar cell (green), presynaptic amacrine cell (cyan). (\textbf{F}) Summary of the number of contacts color-coded by the cell type of the contacting cell: AII (blue), cone bipolar (green), presynaptic amacrine cell (cyan), other (black). Scale bars: 2 μm (\textbf{A,E}), 1 μm (\textbf{B–D}). EM: Electron microscopy.

\noindent DOI: 10.7554/eLife.08206.018

\noindent The following source data is available for figure 4:

\noindent\textbf{Source data 1.} Example image stack of an A2 to A2 tight contact.

\noindent DOI: 10.7554/eLife.08206.019

\noindent\textbf{Source data 2.} Example image stack of a cone bipolar to A2 tight contact.

\noindent DOI: 10.7554/eLife.08206.020

\noindent\textbf{Source data 3.} Example image stack of a chemical synapse to A2 cleft contact.

\noindent DOI: 10.7554/eLife.08206.021
\alttext{Same image as Figure 4 of the article; see the description and \texttt{figures/fig4f.csv} in \texttt{main.tex}.}

\pdfpage{12}
\begin{center}
\includegraphics[width=0.5\linewidth,height=0.4\textheight,keepaspectratio]{figures/sfig4-1.png}
\end{center}
\noindent\textbf{Figure 4—figure supplement 1.} Simple measures of contact geometry do not predict gap junctions. (\textbf{A}) Surface area and the volume of the convex hull bounding each contact are plotted for cleft contacts (gray), AII amacrine cell tight contacts (blue) and cone bipolar cell tight contacts (green).

\noindent DOI: 10.7554/eLife.08206.022
\alttext{(A, the only panel) Log–log scatter plot: x ``Surface area (voxels)'' 10\textsuperscript{1}–10\textsuperscript{5}, y ``Convex hull volume (voxels)'' 10\textsuperscript{1}–10\textsuperscript{5}. Grey × marks (cleft contacts) form a narrow band along a straight line from ≈ (80, 160) to ≈ (5 × 10\textsuperscript{3}, 4.5 × 10\textsuperscript{4}), with one outlier low at ≈ (140, 80). Green filled dots (cone bipolar tight contacts) from ≈ (85, 210) to ≈ (3.8 × 10\textsuperscript{3}, 2.4 × 10\textsuperscript{4}) and blue filled dots (AII tight contacts; three visible, two of them touching, at ≈ (400–470, 1.1–1.3 × 10\textsuperscript{3}) and one at ≈ (1.8 × 10\textsuperscript{3}, 8 × 10\textsuperscript{3})) lie within the grey band, not separated from the cleft contacts.}
\figuredata{A}{figures/sfig4-1a.csv}{figures/sfig4-1a.py}{tier 3, digitized from the eLife JPEG (template matching; 3 of 4 blue, 16 of 17 green and 79 of 150 grey markers resolved, the rest hidden by overlap)}

\pdfpage{13}
\begin{center}
\includegraphics[width=0.95\linewidth,height=0.55\textheight,keepaspectratio]{figures/fig5.png}
\end{center}
\noindent\textbf{Figure 5.} Extracellular space preservation increases access to antibodies. (\textbf{A}) Flat-mounted retinas were immunolabeled for vACHT with and without ECS preservation. Retinas were then cross-sectioned and confocal images were taken to directly visualize penetration depth. (\textbf{B}) Same as \textbf{A}, but immunolabeled for vGAT. (\textbf{C}) Fullwidth half-maximal penetration distances (mean +/- SEM) for retinas with and without ECS preservation (n = 20 retina slices from two animals, unpaired Student’s t-test, vACHT: ***p<0.0001, vGAT: ***p<0.0001). (\textbf{D}) EM images collected from the regions of interest in panel \textbf{B} indicated by dotted rectangles. Scale bars: 50 μm (\textbf{A,B}), 1 μm (\textbf{D}).

\noindent DOI: 10.7554/eLife.08206.023
\alttext{Same image as Figure 5 of the article; see the description and \texttt{figures/fig5c.csv} in \texttt{main.tex}.}

\end{document}
